\section{Introduction}

Single-cell RNA sequencing (scRNA-seq) enables transcriptional profiling at
single-cell resolution, allowing perturbation responses to be resolved that
would be obscured by bulk averaging \citep{wang2021scgnn}.
Genetic perturbation designs that combine knockout or overexpression with
single-cell readouts are now routine, yet their analysis typically relies on
principal component analysis (PCA) followed by clustering---an approach that
treats the data as a homogeneous point cloud and discards the typed relational
structure among cells, the genes they express, and the regulatory modules
those genes form \citep{chervov2022cellcycle}.
In the MCF10A epithelial cell system examined here, p53 status crossed with
three CENP-A overexpression states defines six biologically distinct perturbation
conditions \citep{jeffery2022cenpa}.
Because the perturbation effects are large, classical linear classifiers approach
ceiling performance, creating conditions where differences in representation quality
are masked by near-identical classification accuracy.

Graph neural network (GNN) methods have emerged as principled alternatives for
single-cell analysis because they can learn over non-Euclidean cell-graph
topologies.
scGNN demonstrated that graph attention over cell neighborhoods yields richer
representations for trajectory inference and imputation than k-means
\citep{wang2021scgnn}.
DeepMAPS used a heterogeneous graph transformer to construct cell-type-specific
gene regulatory networks from paired scRNA and ATAC data, showing that typed
graph structure captures regulatory information invisible to matrix factorization
\citep{ma2023deepmaps}.
MarsGT extended this to multi-omics settings, identifying rare cell populations
from heterogeneous graphs of cells, proteins, and chromatin accessibility nodes
\citep{wu2024marsgt}.
Perturbation-specific frameworks such as GEARS leverage gene interaction graphs
to predict multi-gene perturbation outcomes, highlighting how graph structure over
gene relationships recovers combinatorial perturbation effects that linear models
miss \citep{roohani2023gears}.
Probabilistic latent-variable models such as scVI provide competitive latent
representations when large amounts of labeled data are available, but optimize
for reconstruction and batch correction rather than perturbation-discriminative
cell embeddings \citep{lopez2018scvi}.

Heterogeneous Graph Transformers (HGT) \citep{hu2020hgt} generalize standard GNN
message passing by assigning type-specific projection matrices to each
(source type, edge type, destination type) triplet.
This typed parameterization avoids the collapse problem of homogeneous methods:
cell--cell similarity, cell--gene expression, gene--gene coexpression, and
gene--module membership are each modeled with dedicated attention weights.
When applied to single-cell perturbation data, this structure aligns naturally
with the biological reality that p53 loss and CENP-A overexpression remodel
cell neighborhoods, perturb gene coexpression modules, and alter gene-to-module
membership simultaneously.

This study presents a reproducible HGT pipeline for the MCF10A p53/CENP-A
dataset (E-MTAB-9861), benchmarks it under five-fold stratified cross-validation
against k-means, logistic regression, random forest, MLP, and a homogeneous GCN
ablation, and characterizes learned cell embeddings using UMAP visualization
\citep{mcinnes2018umap}, silhouette score \citep{rousseeuw1987silhouette},
Calinski-Harabasz index \citep{calinski1974dendrite}, and label-efficiency
analysis.
The main finding is that near-ceiling linear classification accuracy masks
substantial differences in representation quality: HGT-derived embeddings achieve
silhouette scores 40-fold higher than raw PCA coordinates, producing
biologically interpretable representations that jointly encode p53 and CENP-A
perturbation structure.
